\cotefile{79}
\scanpage{250}

\unsure{demment|évidemment} surjectif, dont le noyau est le groupe

\[
\unsure{\mathfrak{z}|\mathfrak{Z}}=\deleted{\unsure{G_0}}G_{(s_0,t_0)}=\{g\in G\mid g(s_0)=s_0,\ g(t_0)=t_0\}\tag*{(91)}
\]

% slanted
\margin{NB. On a $\mathfrak{z}=G_0(s_0,t_0)$}

d'où on a une suite exacte

\[\begin{aligned}
1&\unsure{-|\longrightarrow}\mathfrak{z}\unsure{-|\longrightarrow}T_G\longrightarrow Z\times Z\unsure{-|\longrightarrow}1\\
&\hspace{5em}g\longmapsto(\delta_0(g),\delta_1(g))
\end{aligned}\tag*{(92)}\]

On pose

\[
\left\{\begin{aligned}
T_{0\deleted{G}}&=\operatorname{Ker}\delta_0=\deleted{\unsure{G_{s_0}}\unsure{\cap T_G}=}\{g\in G\mid gs_0=s_0,\,gt_0\in Zt_0\}\\
T_{1\deleted{G}}&=\operatorname{Ker}\delta_1=\{g\in G\mid gs_0\in Zs_0,\,gt_0=t_0\}
\end{aligned}\right.\tag*{(93)}
\]

Notons qu'on a

\[
\left\{\begin{aligned}
&\deleted{\unsure{T_G}}\mathfrak{z}\subset T_{0\deleted{G}},T_{1\deleted{G}}\subset\deleted{\unsure{G_0}}T_G,\quad Z\subset T_G\\
&T_G=T_{0\deleted{G}}\cdot T_{1\deleted{G}}=T_{0\deleted{G}}\cdot Z=\deleted{Z}T_{1\deleted{G}}\cdot\deleted{Z}\\
&T_{0\deleted{G}}\cap T_{1\deleted{G}}=T_{0\deleted{G}}\cap Z=Z\cap T_{1\deleted{G}}=\mathfrak{z}
\end{aligned}\right.\tag*{(94)}
\]

et

\[
\left\{\begin{aligned}
&\sigma(T_G)=T_G\\
&\sigma(T_{0\deleted{G}})=T_{1\deleted{G}},\quad \sigma(T_{1\deleted{G}})=T_{0\deleted{G}}\\
&\text{et bien sûr}\qquad \sigma(Z)=Z\quad (Z\text{ central }?)
\end{aligned}\right.\ .\tag*{(95)}
\]

On a aussi

\[\begin{aligned}
T_G\cap H=T_H&\subset\operatorname{Ker}(\delta_0.\delta_1:T_G\longrightarrow Z)\\
&\hspace{9em}\deleted{\unsure{g}}t\longmapsto\delta_0(t)\delta_1(t)
\end{aligned}\tag*{(96)}\]

et \deleted{\unsure{sur}} $\delta_H:T_H\xrightarrow{\sim}Z$ est \unsure{égale à la|opposée à la} restriction de $\unsure{\delta_0}$ à $T_H$

\deleted{Introduisons maintenant}

\[\deleted{T_G\ \text{normalise }U_0\text{ et centralise }T_H}\tag*{\deleted{(97)}}\]

\[\deleted{\operatorname{Ker}(\delta_0\delta_1:T_G\longrightarrow Z)=T_H}\tag*{\deleted{(97)}}\]

\deleted{Comme \unsure{$G$} $T_{G\unsure{0,1}}$ est défini comme le centralisateur}

La formule (96) implique que $T_G$ normalise $T_H$ (puisque $H$ est distingué dans $G$), donc comme $T_G/\mathfrak{z}$ est commutatif et $T_H\cap\mathfrak{z}\subset\unsure{\{H\cap\mathfrak{z}\}|H\cap\mathfrak{z}}=\{1\}$, on voit

\[T_G\text{ centralise }T_{\deleted{\unsure{G}}H}\tag*{(97)}\]

% archive sheet number: 249
